\section{Static Perturbation analysis of large-scale binary cell-fate decision networks} \label{ssec:stat_pert}
In this section, we investigate the impact of static perturbations (to edge weights and density; see \S\ref{sssec:noise_def}) on phenotypic distributions, team strengths, and their interrelation. First, we qualitatively understand the effect of noise on the steady-state distribution. We will display results for one artificial and biological network each, and provide reference to the corresponding section in the supplementary material for other networks. Fig.~\ref{fig:S1} provides the frequency of (steady-state) phenotypic distributions of $\emttwosix$ (top) and $\tseight$ (bottom) for three different noise levels (legend; arranged in decreasing order of noise level) of static perturbations to edge weights (left) and density (right). The dots represent the mean of the frequencies of the state across the iterations they have occurred. The upper and lower error bars represent the $75$-th and $25$-th percentile of this distribution, respectively. We characterize the steady states by the difference between the team expressions (see \S\ref{text:team_exp}) of the bigger and the smaller teams, i.e., $\mathrm{T1-T2}$, as shown in the x-axis. The black dashed line belongs to the corresponding wild-type network. 

Our general qualitative observations are: One, the dominant state in the unperturbed network remains dominant in perturbed networks up to $\tsw = 0.17$ for biological networks and $\tsw = 0.61$ for artificial networks, below which the state with $\mathrm{T1-T2}$ in the wild-type network closer to $0$ becomes dominant. In the latter case, prediction of the exact dominant state across noise levels isn't possible as they occur non-uniformly. Two, the number of hybrid states (states with team expressions that are neither $(0,1)$ nor $(1,0)$) increases with an increase in noise level. Three, the states that occurred in unperturbed or relatively less perturbed networks continue to occur in networks with higher perturbations. Four, there is no pattern in the frequencies for a given hybrid state across changes in noise levels. The above observations are common regardless of the network-specific nuances and the differences between the artificial and biological networks. The one contrast between artificial and biological networks is that the dominant states are always the pure states in unperturbed artificial networks, whereas, in biological networks, there are some networks with dominant states that are not pure even for $\tsw$ above the thresholds mentioned above. This is probably because of the higher interconnectivity and smaller size of artificial networks than biological networks. The frequency of dominant states decreases with increasing noise in all cases except for biological networks subject to edge-weight perturbations. Another observation for perturbed networks with wild-type dominant states not being pure is that the frequency of pure states increases with increasing noise at the expense of the dominant state frequency. In conclusion, for networks with team strength above a threshold, the teaming behavior supports the unperturbed low-dimensional phenotypic distribution. In the absence of a 'strong' teaming behavior, the network tries to behave as one unit.

\subsection{Static Edge-weight Perturbations} \label{sssec:stat_ew}
In this section, we quantify our qualitative observations for edge-weight perturbations, as summarized in Fig.~\ref{fig:S2}. We obtain the mean of team strength distributions ($\ts$) from the iterations (see App.~\ref{ap:ts_jsd_dist}) and normalize them by the corresponding team strength of the unperturbed (wild-type) network ($\tsw$). To quantify $\ts/\tsw$ as a function of edge-weight noise levels, we fit the data to an empirical model defined as:
\begin{align}
    \ts(\ewnum) &\equiv \ts, \\
    \ts(\tsw,\ewnum) &= p\cdot\tsw\exp(q\cdot\ewnum) - r, \label{eq:ts_ewnum}
\end{align}
where $\ewnum$ is the numerical value of a given edge-weight range (see Eq.~\ref{eq:rw_ew}; \S\ref{tt:ew_range_num}). The team strength values obtained from the resulting fits are plotted (as black crosses) along with $\ts$ (blue dots) against the corresponding edge-weight ranges, across networks, as shown in the top-left panel in Fig.~\ref{fig:S2}. The edge-weight ranges (x-axis) are arranged in the ascending order of the noise level they introduce, i.e., the $\ewnum$. The black-dashed line is Eq.~\ref{eq:ts_ewnum} calculated over the edge-weight noise range. The fit parameters are shown in Tab.~\ref{tab:ew_param}. We emphasize that the fit parameters aren't independent of each other and hold no physical property. This demonstrates that the teams' resilience to static edge-weight perturbations is attributed to its underlying team-based topology of the network, quantified using team strength. A natural expectation is that the emergent phenotypes from teams must also be resilient to static edge-weight perturbations, with the latter increasing with team strength, and we test this below. 

The middle and bottom panels in Fig.~\ref{fig:S2} correspond to an artificial and a biological network, respectively. We use JSD to quantify the divergence of phenotypic distribution of the perturbed network from the corresponding wild-type network, with greater JSD implying greater divergence. Therefore, higher Jensen-Shannon Divergence ($\jsd$) marks a lower resilience of the network to the introduced noise.  (see \S\ref{{sssec:jsd_def}}). The top-right panel in Fig.~\ref{fig:S2} shows the JSD distributions across edge-weight noise levels. We see a general trend in the variance of the JSD distributions, increasing with increasing edge-weight noise level. We are interested in the role of team strength on JSD in the presence of perturbations. Therefore, we propose an empirical model to establish a relationship between $\jsd$ and $\ts$ for a given network, defined as:
\begin{equation} \label{eq:jsd_ts_net}
    \jsd = a\cdot\log\left(\frac{\tsw-\ts}{\ts}\right)+b,
\end{equation}
where $b=\jsd(\ts=\tsw/2)$. Here, $\jsd$ refers to the mean of the JSD distribution, unless mentioned otherwise. We fit Eq.~\ref{eq:jsd_ts_net} to $\jsd$ and $\ts/\tsw$ across edge-weight ranges (legend) as shown in the middle and bottom-left panel in Fig.~\ref{fig:S2}. The black dashed line interpolates Eq.~\ref{eq:jsd_ts_net} over the range of $\ts/\tsw$, i.e., $[0,1]$. Note that the orange dots at the extremes are not points obtained from our simulation but correspond to $\jsd(\ts = 0) = 1$ and $\jsd(\ts = 1) = 0$ according to Eq.~\ref{eq:jsd_ts_net}. Tab.~\ref{tab:ew_param} displays the obtained parameters $a$ and $b$ from LMFIT along with their 1-sigma variations for each network. As expected, $b$ decreases with an increase in $\tsw$. We see no trend in $a$ as the first term on the right of Eq.~\ref{eq:jsd_ts_net} is independent of $\tsw$ (see Eq.~\ref{eq:ts_ewnum}). Thus, we can conclude that, for a given network, the resilience of the phenotypic distributions to static edge-weight perturbations increases with increasing team strength.

To check if the models Eqs.~\ref{eq:ts_ewnum} and~\ref{eq:jsd_ts_net} can predict the dependence of $\jsd$ on $\ts$ across networks for a particular edge-weight range, Eq.~\ref{eq:ts_ewnum} is plugged into Eq.~\ref{eq:jsd_ts_net} to result the model:
\begin{equation} \label{eq:jsd_ts_ran}
    \jsd = a\cdot\log\left(\frac{1}{p\cdot\exp(q\cdot\ewnum) - r}-1\right)+b.
\end{equation}
The $\jsd$-values along with their error-bars are calculated for constant $\ewnum$s using parameter values from Tabs.~\ref{tab:ew_art_all} and~\ref{tab:ew_art_all}. These are displayed alongside $\jsd$ (obtained from our simulations) plotted against $\ts/\tsw$ for a particular edge-weight range in the middle and bottom-right panel in Fig.~\ref{fig:S2}. We emphasize that no further fitting is done here, and Eq.~\ref{eq:jsd_ts_ran} predicts our data within 2-sigma variations. While $\jsd$ decreases with an increase in $\tsw$ for artificial networks as expected, the trend is the opposite for biological networks. This observation remains unexplained.

\section{Methods}

\subsection{Jensen--Shannon divergence}\label{sssec:jsd_def}
Jensen--Shannon divergence (JSD) is a symmetric, bounded measure of similarity between two probability distributions \cite{menendez1997jensen}, derived from the KL divergence:
\begin{align}
    \jsd(P \| Q) &= \frac{\mathrm{KL}(P \| M) + \mathrm{KL}(Q \| M)}{2}, \quad M = \frac{P + Q}{2}, \label{eq:jsd_def} \\
    \mathrm{KL}(P \| M) &= \sum_{x \in X} P(x) \log_2 \frac{P(x)}{M(x)}.
\end{align}
JSD takes values in $[0, 1]$, with $\mathrm{JSD} = 0$ indicating identical distributions. It is used here to quantify the divergence of the perturbed network's phenotypic distribution from the wild-type, serving as a quantitative measure of network resilience: higher JSD implies lower resilience.

\section{Fitting and error propagation}\label{ssec:fit_meth}
LMFIT \cite{https://doi.org/10.5281/zenodo.598352} is used for model fitting. Error bars ($1\sigma$ intervals) are calculated using Monte Carlo error propagation, which does not require the assumption that the distributions under consideration are Gaussian. Parameter values obtained from LMFIT are treated as Gaussian, and random samples drawn from these distributions are propagated through the model to generate a distribution for each derived quantity. The mean, median, and $16$th and $84$th percentiles are extracted from the resulting distribution; the upper and lower error limits are defined as the $84$th percentile minus the median and the median minus the $16$th percentile, respectively.

\subsection{Static edge weight perturbation}

To complement the continuous noise simulations and characterise the steady-state sensitivity of network phenotypes to fixed changes in interaction strength, we also performed static perturbation experiments. For each simulation run, a weight for each edge was independently sampled from a uniform distribution within one of the following ranges: $[0, 0.25]$, $[0, 0.5]$, $[0, 0.75]$, $[0, 1]$, $[0.25, 1]$, $[0.5, 1]$, $[0.75, 1]$, and $[0.95, 1]$, ordered from highest to lowest noise level. The topology and sign of interactions were held constant across all perturbations.

\noindent\textit{Assigning numerical values to edge-weight ranges:} As the smallest interval considered is $0.05$, the range $[0, 0.05]$ is assigned value $1$. The numerical value for an arbitrary range is defined as:
\begin{align}
    \ewnum([0, \mathrm{up}_{lim}]) &= \frac{\sum_{i=1}^{j} i}{j}, \quad j = \frac{\mathrm{up}_{lim}}{0.05} \nonumber \\
    \ewnum([\mathrm{low}_{lim}, 1]) &= \frac{\ewnum([0,1]) \cdot 20 - \ewnum([0, \mathrm{low}_{lim}]) \cdot k}{20 - k}, \quad k = \frac{\mathrm{low}_{lim}}{0.05} \nonumber \\
    \ewnum([\mathrm{low}_{lim}, \mathrm{up}_{lim}]) &= \frac{\ewnum([0, \mathrm{up}_{lim}]) \cdot j - \ewnum([0, \mathrm{low}_{lim}]) \cdot k}{20 - k}. \label{eq:rw_ew}
\end{align}
Note that this metric is undefined when $\mathrm{low}_{lim} = \mathrm{up}_{lim}$.


\paragraph{Initial conditions} For networks with $N < 20$, all $2^N$ possible initial states are enumerated. For $N \geq 20$, $2^{20}$ randomly generated initial states are used. In both cases, an unperturbed (wild-type) simulation with edge-weight moduli of $1$ is run as a baseline against which phenotypic distributions under perturbation are compared.
